\section{Implantation}
\label{sec:reevaluation-base-regles-implantation}

Les procédures permettant le support de la suppression d'orbite et le
déplacement d'une opération suite à l'édition d'une spécification paramétrique
sont, au moment de la rédaction de ce manuscrit, des travaux en cours. %
Dans cette section, nous détaillons les implantations des structures de données
des DAGs d'évaluation et de réévaluation grâce auxquels nous pouvons apparier
des orbites. %
Pour permettre le mécanisme de réévaluation que nous proposons, nous avons
implanté les structures illustrées dans les
figures~\ref{fig:reevaluation-scripts-implantation-struct-dag-eval} et
\ref{fig:reevaluation-scripts-implantation-struct-dag-reeval}. %

\begin{figure}[p!]
	\centering
	\includegraphics[width=.8\textwidth]{class-diagram-hr-datastructure}
	\caption{Diagramme de la classe \textit{Evaluation\-DAG}}
	\label{fig:reevaluation-scripts-implantation-struct-dag-eval}
\end{figure}

\begin{figure}[p!]
	\centering
	\includegraphics[width=.8\textwidth]{class-diagram-rt-datastructure}
	\caption{Diagramme de la classe \textit{Reevaluation\-DAG}}
	\label{fig:reevaluation-scripts-implantation-struct-dag-reeval}
\end{figure}

Les DAGs d'évaluation et de réévaluation sont deux structures représentant
l'historique d'évolution d'une orbite dans deux contextes différents~:
l'évaluation et la réévaluation. %
Ainsi, un DAG trace un historique au travers des opérations enregistrées dans
une spécification paramétrique. %

Pour l'implantation de ces structures, nous avons fait le choix de représenter
une application en deux classes distinctes qui décrivent un niveau
d'application. %
La première classe \textit{EventLevel} décrit une application du point de vue de
la règle. %
Cette classe spécifie qu'un niveau d'évènements détient un numéro d'application
\textit{appNum}, un type d'application \textit{appType} et la liste des
évènements détectés \textit{eventList} pour une orbite donnée. %
Notons que lors de l'évaluation initiale, \textit{appType} est initialisée à
\textit{INIT}. %
La deuxième classe \textit{OrbitLevel} décrit une application du point de vue
d'une orbite suivie. %
Dû à la différence de contexte dans lequel nous utilisons les DAGs d'évaluation
et de réévaluation, ces deux classes sont implantées comme des classes
abstraites. %
Une instance \textit{EventLevelEval} (ou \textit{EventLevelReeval}) détient
également une référence vers le niveau d'orbite relatif à la même application. %
Pour les niveaux d'orbites, rappelons que si à l'évaluation nous enregistrons
des noms n\oe uds afin d'être le plus générique possible, ce n'est pas le cas
pour la réévaluation. %
En effet, dans la phase d'appariement, un DAGs de réévaluation enregistre
directement des identifiants de brins. %
Ainsi, une instance de \textit{OrbitLevelEval} détient un nom de n\oe ud
\textit{nodeName} et une référence vers le niveau d'évènement de l'application
suivante. %
Au contraire, une instance de \textit{OrbitLevelReeval} détient un identifiant
de brin \textit{dartID} et une référence vers le niveau d'évènement de
l'application suivante. %

À l'intérieur de ces niveaux d'applications, nous organisons en alternance des
orbites et des évènements et leurs relations de suivi ou d'origine. %
Pour représenter les d'évènements, nous implantons la classe \textit{EventNode}
qui détient un tag de type \textit{Event} (une simple énumération des évènements
que nous pouvons détecter), un numéro de branche \textit{branchIndex} et une
référence vers une unique orbite \textit{child}. %
Pour les orbites, nous implémentons la classe \textit{OrbitNode} qui détient un
type d'orbite \textit{orbitType}, un chemin \textit{alphaPath}, un numéro de
branche \textit{branchIndex} et une liste d'enfants \textit{children}. %
Avec cette implantation, nous exprimons deux faits, un n\oe ud d'évènement ne
peut référencer qu'un seul n\oe ud d'orbite, car tout évènement dans une règle
n'est étudié que pour une orbite donnée. %
Au contraire, une orbite peut être impactée par plusieurs évènements. %
Notons également que l'enfant d'un n\oe ud d'orbite est une instance de la
classe \textit{Link} que nous créons pour associer un n\oe ud d'évènements à un
type de relation \textit{LinkType}, c'est-à-dire soit un suivi \textit{TRACE},
soit une origine \textit{ORIGIN}. %

Enfin, nous avons décidé d'encapsuler les niveaux d'application dans une liste
afin de pouvoir itérer le contenu d'un DAG. %
Toutefois, cela ne considère pas la possibilité que plusieurs branches soient
nécessaires pour représenter pleinement l'historique des évolutions d'une
orbite. %
Ainsi, pour les DAGs d'évaluation, nous décidons d'encapsuler la structure de
données dans une \textit{Map} où chaque niveau d'application est accessible
directement à partir de son \textit{appID}. %
Le choix d'une \textit{Map} revêt un aspect pratique, car au cours de la
procédure de réévaluation nous avons besoin de récupérer des morceaux de DAGs
d'évaluation niveau par niveau. %
Au contraire, pour les DAGs de réévaluation, la structure de données est
simplement encapsulée dans une autre liste, puisque l'accès se fait toujours par
rapport au dernier niveau inséré. %


%%% Local Variables:
%%% mode: LaTeX
%%% TeX-master: "../../../main"
%%% End:
